\section{Maîtrise du coût de compilation du Valo Engine}
\label{chap:t2}
% Entreprise : Free Mobile — Chapitre « outils : productivité et
% standardisation de l'environnement » (section 2/3)

% =====================================================================
%  Trame conseillée (réutilisée pour chaque tâche) :
%  1. Contexte & dette technique de départ
%  2. Objectif de la mission
%  3. Conception / choix techniques
%  4. Réalisation (extraits de code COMMENTÉS, schémas — PAS de code entier)
%  5. Difficultés rencontrées
%  6. Résultats & bénéfice business (réduction de la dette)
% =====================================================================

\subsection{Contexte et dette technique initiale}
Le \emph{Valo Engine} est écrit en \textbf{C++}, langage dont la compilation est
réputée \textbf{coûteuse}, en temps de calcul comme en mémoire. Cette lourdeur
tient à son modèle historique de compilation, hérité du C : chaque fichier source
est compilé en une \textbf{unité de traduction} indépendante, qui inclut
\emph{textuellement} les en-têtes (\emph{headers}) dont elle dépend. Les
\textbf{templates} aggravent ce mécanisme : leur définition résidant dans les
en-têtes, ils sont \textbf{réinstanciés à chaque unité de traduction} qui les
inclut. Une même instanciation est ainsi recompilée des dizaines de fois à
travers le projet.

À strictement parler, ce coût ne devient un problème que parce que l'on
cherche à compiler \textbf{en parallèle} : une compilation séquentielle
resterait sobre en mémoire, mais à un prix prohibitif en temps. Or y
renoncer n'est pas une option : développer en C++, c'est s'appuyer sur le
compilateur pour se prémunir d'un maximum d'erreurs \emph{avant même}
d'exécuter le programme, et la rapidité du \emph{build} conditionne donc
toute la boucle de développement (lancer les tests unitaires, valider une
modification), avec les gains de temps évoqués en introduction de ce
chapitre.

Le pic se produit à l'\textbf{édition des liens}. L'éditeur de liens doit y
rassembler tous les fichiers objets et \textbf{dédupliquer} les définitions
redondantes d'une même fonction ou instanciation, issues des multiples unités de
traduction ; l'empreinte mémoire de cette étape explose à mesure que le nombre de
symboles croît. En lançant une compilation parallèle exploitant l'ensemble des
cœurs d'une machine, de l'ordre de \textbf{24}, la \gls{toolchain} \textbf{saturait
la RAM} précisément à cette phase de liaison, rendant la compilation instable
voire impossible à mener à son terme. La dette était ici nichée dans la
\textbf{chaîne de compilation} elle-même : un atout de productivité (le
parallélisme) se retournait en obstacle.

\subsection{Objectif de la mission}
La mission consistait à étudier la \textbf{faisabilité d'un compromis} : exploiter
au mieux les cœurs disponibles \emph{sans} provoquer la saturation mémoire qui
faisait échouer la compilation. Autrement dit, \textbf{ordonnancer} la
compilation pour qu'elle reste \textbf{soutenable} en mémoire, en particulier
durant l'édition des liens, tout en conservant un parallélisme aussi élevé que
le permet le budget RAM de la machine.

\subsection{Conception et choix techniques}
Le diagnostic a d'abord consisté à \textbf{isoler la phase critique}. La
compilation des unités de traduction se parallélise sans difficulté majeure ;
c'est l'\textbf{édition des liens}, beaucoup plus gourmande en mémoire par
processus, qui fait diverger la consommation lorsque plusieurs liaisons
s'exécutent simultanément. La piste retenue revient donc à \textbf{découpler le
degré de parallélisme} des deux phases : conserver un parallélisme large pour la
compilation, mais \textbf{borner} le nombre de tâches d'édition des liens
exécutées de front, afin de plafonner le pic mémoire sous la capacité réelle de
la machine.

% TODO : préciser le mécanisme retenu (pool de tâches de l'outil de build,
%        limitation dédiée à l'étape de link, choix de l'éditeur de liens, etc.).
À plus long terme, la \textbf{cause racine}, à savoir la recompilation répétée
des en-têtes et des templates, trouve une réponse dans les \textbf{modules} de
\textbf{C++20}, qui compilent une interface \textbf{une seule fois} au lieu de la
réanalyser à chaque inclusion. Cette voie n'était toutefois pas exploitable :
\textbf{le support des modules reste partiel} dans les compilateurs. À titre de
repère, sur Ubuntu~24.04 LTS, \texttt{g++}~13.3 conserve \textbf{C++17} comme
standard \emph{par défaut} (\verb|__cplusplus = 201703L|) ; l'adoption de C++20 ou
C++23, et de leurs modules, supposerait une migration de la \gls{toolchain} qui
n'est pas encore une option standard et mûre.
% TODO (Claude) : à la réécriture, revérifier le standard C++ par défaut et le
%   compilateur le plus avancé sur l'environnement de PRODUCTION ciblé (la valeur
%   ci-dessus a été constatée localement sur g++ 13.3 / Ubuntu 24.04).

\subsection{Réalisation}
% TODO (auteur) : insérer l'extrait réel (configuration de build / script
%        d'ordonnancement) si disponible.
L'idée se matérialise en \textbf{limitant la concurrence de l'étape de liaison}
indépendamment de celle de la compilation, de sorte que le pic mémoire cumulé
des éditions de liens reste sous le plafond de la machine.

\subsection{Difficultés rencontrées}
La principale difficulté a été de \textbf{caractériser le pic mémoire} : la
saturation ne survenait qu'à l'édition des liens et seulement sous forte
concurrence, ce qui en rendait l'observation peu intuitive : le coût n'est pas
réparti uniformément sur la durée de la compilation, mais \textbf{concentré} sur
sa phase finale. La seconde tenait à la \textbf{tension entre vitesse et
mémoire} : tout plafonnement de la concurrence de liaison améliore la stabilité
au prix d'un allongement du temps de compilation, et le bon réglage dépend
étroitement du couple (mémoire disponible, taille des objets à lier).

\subsection{Résultats et bénéfice business}
% TODO : quantifier (compilation menée à terme sans crash, temps obtenu vs
%        configuration naïve). Relier au thème dette technique (chap.~\ref{chap:veille}).
En bornant la concurrence de la phase de liaison, la compilation redevient
\textbf{soutenable} sur les machines de développement : elle peut être menée à
son terme sans saturer la RAM, là où une parallélisation naïve la faisait
échouer, tout en conservant un parallélisme maximal sur la phase de
compilation proprement dite, qui représente l'essentiel de la durée totale.
Pour l'équipe, le bénéfice est direct : un \emph{build} fiable est la
condition de toute boucle de développement courte, et donc de toute
intervention sereine sur le code.

Au-delà du résultat immédiat, cette mission a été l'occasion de
\textbf{mettre au jour une dette de fond}, le coût du modèle de compilation C++
préalable aux modules, dont la résorption complète reste suspendue à la
maturation du support de C++20/23 dans les compilateurs.
% TODO (auteur) : si disponible, chiffrer le temps de compilation obtenu et
%        le seuil de concurrence retenu pour l'édition des liens.
